%-------------------------------------------------------------------------------
%	SECTION TITLE
%-------------------------------------------------------------------------------
\cvsection{个人简介}


%-------------------------------------------------------------------------------
%	CONTENT
%-------------------------------------------------------------------------------
\begin{cvparagraph}

Linux 系统构建与服务交付方向，覆盖多架构交叉编译、内核维护、镜像交付与容器化服务部署。\\
技术栈按工作场景分组：\\
\textbf{系统构建}：OpenWrt、Armbian、交叉编译、内核配置、Device\nobreakspace{}Tree、镜像打包、UEFI、RISC-V64 / ARM64 / MIPS / x86\\
\textbf{内核与调试}：C、QEMU / GDB、串口与内核日志\\
\textbf{容器与服务交付}：Docker、docker-compose、服务启动、健康检查、日志轮转与异常恢复\\
\textbf{网络与终端运维}：TCP/IP、VLAN / VPN / 防火墙、LNMP（Nginx / MySQL / PHP）、办公终端与业务账号权限\\
\textbf{本地 AI 服务}：Ollama、模型量化、LoRA、CUDA / AVX\\
\textbf{脚本与工具}：Shell / Python、Go（基础）、Make / CMake / Ninja、Git\\
维护过 RISC-V / ARM / MIPS / x86 平台的 \mbox{OpenWrt}、Armbian 构建环境，使用 QEMU / GDB 定位内核启动和驱动问题，并有 \mbox{OpenWrt} 上游补丁合入经历。\\
在边缘设备上完成 Docker 化 LLM 服务的部署与运行维护；处理过 VLAN、VPN、防火墙、业务账号权限及办公终端问题。\\
习惯使用 Shell / Python 编写构建、服务检查、日志采集与恢复脚本；使用 Hermes、DeepSeek Harness 辅助定位问题并人工验证。
\end{cvparagraph}
